\clearpage
\section{Quantization of the Field}
\begin{colbox}
	\begin{enumerate}
		\item Calculate the commutators \([v_x,H]\), \([v_y, H]\), \([v_z,H]\)
		\item Use Ehrenfest's theorem to calculate \begin{align}
			C_1(t) = \left\langle S_zv_z \right\rangle && C_2(t) = \left\langle S_xv_x+S_yv_y \right\rangle && C_3(t) = \left\langle S_xv_y - S_yv_x \right\rangle
		\end{align}
		\item Write the time development of \(\left\langle\vec{S}\cdot\vec{v}\right\rangle\) as a function of \(C_1(0), C_2(0), C_3(0)\)
		\item a beam of electrons moving at velocity \(\vec{v}\) is prepared at \(t=0\). The electrons in the beam have spin such that \(C_1(0), C_2(0), C_3(0)\) are known. The electrons in the beam interact with the magnetic field \(\vec{B}\) in the time frame \([0,T]\). At time \(t=T\) a value corresponding to \(\left\langle\vec{S}\cdot\vec{v}\right\rangle\). Conclude \(a\) from the graph.
		\item Compare the experiment result to theory.
	\end{enumerate}
\end{colbox}

\subsection{Commutators}
We'll represent both all the operators using the \(\Pi\) operators as we did in the lecture. If we define
\begin{equation}
	\Pi_i = \frac{p_i-aA_i}{\sqrt{\hbar eB}}
\end{equation}
then
\begin{align}
	H = \frac{\hbar \omega_c}{2}\ins{\Pi_x^2+\Pi_y^2+\Pi_z^2-\sigma_z\ins{1+a}}
\end{align}
and
\begin{equation}
	\vec{v} = \frac{\sqrt{\hbar eB}}{m}\vec{\Pi}
\end{equation}
therefore
\begin{align}
	[v_x,H] &= \frac{\sqrt{\hbar eB}}{m}\frac{\hbar\omega_c}{2}[\Pi_x,\Pi_y^2] = i\hbar\omega_c v_y \\
	[v_y,H] &= \frac{\sqrt{\hbar eB}}{m}\frac{\hbar\omega_c}{2}[\Pi_y,\Pi_x^2] \\
	[v_z,H] &= 0
\end{align}
\subsection{Ehrenfest}
To prepare I'll calculate the commutators
\begin{align}
	[S_x,H] &= -\frac{\hbar^2\omega_c}{4}\ins{1+a}[\sigma_x,\sigma_z] = i\hbar\omega_c\ins{1+a}S_y \\
	[S_y,H] &= -\frac{\hbar^2\omega_c}{4}\ins{1+a}[\sigma_y,\sigma_z] = -i\hbar\omega_c\ins{1+a}S_x \\
	[S_z,H] &\propto [\sigma_z,\sigma_z] = 0
\end{align}
and thus
\begin{align}
	\diff*{C_1}{t} &= \frac{1}{i\hbar}\langle[S_z v_z, H]\rangle = 0 \\
	\diff*{C_2}{t} &= \frac{1}{i\hbar}\langle[S_x v_x + S_y v_y, H]\rangle = 0 \\
	&= \frac{1}{i\hbar}\langle i\hbar\omega_c\ins{1+a}S_y v_x +i\hbar\omega_c v_y S_x - i\hbar\omega_c\ins{1+a}S_x v_y - i\hbar\omega_c S_y v_x \rangle \\
	&= \Omega \langle S_y v_x + S_x v_y \rangle\\
	&= -\Omega C_3 \\
	\diff*{C_3}{t} &= \frac{1}{i\hbar}\langle[S_x v_y - S_y v_x, H]\rangle \\
	&= \frac{1}{i\hbar}\langle i\hbar\omega_c\ins{1+a}S_y v_y - i\hbar\omega_c v_x S_x + i\hbar\omega_c\ins{1+a}S_x v_x - i\hbar\omega_c S_y v_y \rangle \\
	&= \Omega \langle S_y v_y + S_x v_x \rangle\\
	&= \Omega C_2 \\
\end{align}
in total we have
\begin{align}
	\diff*{C_1}{t} &= 0\\
	\diff*{C_2}{t} &= -\Omega C_3\\
	\diff*{C_3}{t} &= \Omega C_2
\end{align}
differentiating the equation for \(C_3\)
\begin{equation}
	\diff*[2]{C_3}{t} = \Omega\diff*{C_2}{t} = -\Omega^2 C_3 
\end{equation}
whose solutions are a linear combination of sines and cosines. therefore
\begin{align}
	C_1 &= A \\
	C_2 &= D\sin\Omega t + E\cos\Omega t \\
	C_3 &= E\sin\Omega t - D\cos\Omega t
\end{align}
since \(C_2, C_3\) have opposite signs in the first order and depend on each other.
\subsection{S dot v}
\begin{align}
	\langle\vec{S}\cdot\vec{v}\rangle &= \langle S_x v_x + S_y v_y + S_z v_z \rangle \\
	&= C_1(t) + C_2(t) \\
	&= C_1(0) + C_2(0)\cos\Omega t - C_3(0)\sin\Omega t
\end{align}
\subsection{Experimental a}
From the graph we can see the period is approximately \(\qty{3.5}{\micro\second}\) therefore
\begin{equation}
	a = \frac{2\pi}{\omega_C T} = \frac{2\pi m}{eBT} = \frac{2\pi \cdot \qty{9.11e-31}{}}{\qty{1.6e-19}{}\cdot\qty{9.4e-3}{}\cdot\qty{3.5e-6}{}} = \sim \qty{0.001}{}
\end{equation}
\subsection{comparison}
The theoretical value is
\begin{equation}
	a = 0.001195
\end{equation}
so the experimental value is very close.